\section{Conclusion}
\label{sec:conclusion}
Graph representation can change which applications enter documented review scope under the same account budget. At five accounts, the YunoHost comparison confirms three name-only and two version-only hits. The follow-ups resolve 164 of 215 pairs and bound turnover at 2.3--26.0\%. In Awesome-Selfhosted, they resolve 88 of 155 pairs, with no confirmed switches and turnover bounded at 0--43.2\%. The direction of net change remains unresolved in both frames. At 50 accounts, the exploratory YunoHost comparison identifies a separate switch.

The graph comparisons quantify a separate tradeoff. At 100 accounts on the original graph, greedy reaches 5,538 repositories and exact union reaches 4,407, with modeled exposure reductions of 11.48\% and 33.79\%, respectively. The coverage--exposure ordering repeats in the separately collected cohort, while the decrease in account--repository pair mass under version continuity changes from 42.44\% to 3.33\%. The exclusion heuristic has a worst deficit of 66 endpoints in the restricted exact benchmarks, and the relative performance of its variants changes across the two graphs.

A fixed-budget comparison can therefore report confirmed additions, removals, and unresolved cases alongside the total number of applications reached. The two consumer frames concern declared production paths at pinned commits and maintainer associations observed at collection time. Evaluating whether a control improves security requires separate evidence on authorization, deployment and intervention outcomes.
